\chapter{深度生成模型}
\begin{quote}\small
\textit{本章摘要。}本章围绕一个问题展开：如何让模型生成符合数据规律、又能用于具体任务的新样本？先区分采样、补全、去噪和条件生成，再以变分自编码器（VAE）、生成对抗网络（GAN）和扩散模型为里程碑，说明它们的生成机制、训练目标和典型失效方式。最后讨论任务选择、样本验证和工业使用边界。
\end{quote}

生成模型不能替真实观测，也不能凭空证明罕见故障确实存在。它可以在数据稀缺时产生候选样本，补全缺失观测，或在给定工况下模拟可能的退化轨迹。合成数据是否有用，要在未参与训练和模型选择的真实数据上检验。

\section{生成任务与概率建模目标}

\subsection{先确定要生成什么}

假设$x\in\R^D$是一个振动窗口，$c$是转速、负载或故障类别等条件。深度生成模型常完成四类任务：
\begin{enumerate}
  \item \textbf{无条件采样：}从学习到的分布中产生新的$x$，例如模拟正常振动窗口。
  \item \textbf{条件生成：}给定$c$产生$x$，例如在指定转速下生成可能的波形。条件通常允许多个结果，不要求给定条件后输出唯一样本。条件生成的目标是近似$p(x\mid c)$，而不是保证每个样本都满足未写入条件变量的隐含约束。
  \item \textbf{补全与去噪：}给定部分观测、带噪观测或掩码，生成与已知信息一致的完整结果。
  \item \textbf{表示学习：}把高维样本压缩为较低维的潜变量，供可视化、聚类或下游预测使用。
\end{enumerate}

这些任务都涉及数据分布，但具体目标不同。无条件生成学习$p_{\theta}(x)$，条件生成学习$p_{\theta}(x\mid c)$，随后可从相应分布采样。补全任务还需把已观测部分作为条件，并保证生成结果与其一致。连续变量的$p_{\theta}(x)$是密度，单个点的密度值不是该点的概率，区域概率要对密度积分。训练集只是有限观测，不能等同真实分布$p_{\mathrm{data}}$。用条件均值替代采样会丢失条件分布中的波动、相关性和多个可能模式。

生成目标与判别目标不同。判别模型通常学习$p(y\mid x)$或直接学习决策边界，生成模型还要描述$x$本身的变化。生成模型的训练损失也不自动等于任务评价指标。重构误差较小，不代表裂纹位置正确；样本看起来逼真，也不代表它覆盖了真实故障类型。

\subsection{潜变量把随机性变成可采样结构}

潜变量$z$是模型引入、但数据中没有直接观测的变量。一个常见生成方向是
\[
z\sim p(z),\qquad x\sim p_{\theta}(x\mid z).
\]
先抽取简单先验中的$z$，再由解码器把它变成样本$x$。例如，若$z\sim\mathcal N(0,1)$且$x\mid z\sim\mathcal N((z,2z)^{\mathsf T},\sigma_x^2I)$，两个观测分量会因共享$z$而相关。观察到$x$后反推$z$，对应后验$p_{\theta}(z\mid x)$。

普通自编码器只要求编码器和解码器互相重构。它没有规定新样本应该从编码空间的哪里抽取。变分自编码器把潜变量先验、概率解码器和近似后验放在同一个目标中，从而得到可采样的表示。

\section{VAE 的潜变量采样}

\subsection{为什么需要变分近似}

变分自编码器（Variational Autoencoder，VAE）用潜变量模型表示联合分布：
\begin{equation}
p_{\theta}(x,z)=p(z)p_{\theta}(x\mid z).
\end{equation}
给定$x$后，真实后验满足
\[
p_{\theta}(z\mid x)=\frac{p_{\theta}(x\mid z)p(z)}{p_{\theta}(x)},\qquad
p_{\theta}(x)=\int p_{\theta}(x\mid z)p(z)\,\mathrm dz.
\]
分母通常无法解析计算。VAE 用编码器给出的$q_{\phi}(z\mid x)$近似该后验，并最大化证据下界（Evidence Lower Bound，ELBO）：
\begin{equation}
\mathcal L(x)=\E_{q_{\phi}(z\mid x)}[\log p_{\theta}(x\mid z)]-\KL(q_{\phi}(z\mid x)\|p(z)).
\label{eq:generative-vae-elbo}
\end{equation}
第一项要求抽到的$z$能够重构$x$。第二项要求不同样本使用的潜变量仍接近先验，以便生成时直接从$p(z)$抽样。

由 Jensen 不等式，
\begin{align}
\log p_{\theta}(x)
&=\log\int q_{\phi}(z\mid x)\frac{p_{\theta}(x,z)}{q_{\phi}(z\mid x)}\,\mathrm dz\\
&\geq \E_{q_{\phi}}[\log p_{\theta}(x,z)-\log q_{\phi}(z\mid x)]
=\mathcal L(x).
\end{align}
更精确的关系是
\begin{equation}
\log p_{\theta}(x)-\mathcal L(x)=\KL\bigl(q_{\phi}(z\mid x)\|p_{\theta}(z\mid x)\bigr)\geq0.
\end{equation}
这里的两个 KL 项用途不同。式\eqref{eq:generative-vae-elbo}中的 KL 比较近似后验与先验，最后一个等式中的 KL 衡量后验近似误差。KL 散度非负但通常不对称，也不是欧氏距离。

\subsection{如何训练和采样}

常用的近似后验为
\[
q_{\phi}(z\mid x)=\mathcal N\bigl(\mu_{\phi}(x),\diag(\sigma_{\phi}^2(x))\bigr).
\]
直接抽样会阻断对编码器参数的普通反向传播。重参数化把随机性移到与参数无关的$\epsilon$：
\begin{equation}
z=\mu_{\phi}(x)+\sigma_{\phi}(x)\odot\epsilon,
\qquad \epsilon\sim\mathcal N(0,I_d).
\end{equation}
因此固定一次$\epsilon$时，梯度可以通过$z$传回编码器。训练阶段给定$x$，从$q_{\phi}(z\mid x)$抽样并计算重构项。生成阶段没有$x$，改为从$p(z)$抽样，再交给解码器$p_{\theta}(x\mid z)$。

若解码器采用$p_{\theta}(x\mid z)=\mathcal N(\mu_{\theta}(z),\sigma_x^2I)$，则
\[
-\log p_{\theta}(x\mid z)=\frac{1}{2\sigma_x^2}\|x-\mu_{\theta}(z)\|_2^2+\frac D2\log(2\pi\sigma_x^2).
\]
所以均方重构项隐含了条件独立、同方差高斯噪声假设。它不一定适合少数但重要的冲击或裂纹像素。

当$q_{\phi}$是对角高斯、$p(z)=\mathcal N(0,I)$时，KL 项为
\[
\KL(q\|p)=\frac12\sum_{j=1}^{d}(\mu_j^2+\sigma_j^2-\log\sigma_j^2-1).
\]
例如一维情况下$\mu=1,\sigma=1$时，KL 为$0.5$。$\mu=0,\sigma=1$时 KL 为零，但这只说明两个分布相同，不说明重构准确。

\subsection{VAE 的变体与典型问题}

VAE 的主要折中是重构准确性与潜变量规整性。KL 权重过强时，编码器可能令$q_{\phi}(z\mid x)$接近先验，潜变量几乎不再携带输入信息，这称为后验坍缩。KL annealing 逐渐增加 KL 权重，free bits 在小 KL 区间内减弱惩罚。它们改变优化折中，却不保证潜变量对应可解释的物理因素。

条件 VAE 将条件加入编码器和解码器，学习$p_{\theta}(x\mid z,c)$。它适合需要连续潜在表示和条件采样的任务。VAE 往往生成较平滑的结果，因为重构目标会偏向条件平均。若细节和少数模式重要，应检查模式覆盖，而不能只增加潜变量维度。

\section{GAN 的对抗训练}

\subsection{判别器如何提供训练信号}

生成对抗网络（Generative Adversarial Network，GAN）把采样器和判别器分开。生成器$G:\R^d\to\R^D$把噪声$z$变成$G(z)$，判别器$D:\R^D\to(0,1)$估计输入来自真实数据的概率。训练时交替更新：判别器区分真实样本与合成样本，生成器则使合成样本更难被判别。采样时只需要生成器。

原始 GAN 目标为
\begin{equation}
\min_G\max_D\;\E_{x\sim p_{\mathrm{data}}}\log D(x)+\E_{z\sim p(z)}\log(1-D(G(z))).
\label{eq:generative-gan}
\end{equation}
固定生成分布$p_g$时，逐点最优判别器为
\[
D^*(x)=\frac{p_{\mathrm{data}}(x)}{p_{\mathrm{data}}(x)+p_g(x)}.
\]
代回式\eqref{eq:generative-gan}，理想情况下目标与
\[
\operatorname{JSD}(p_{\mathrm{data}}\|p_g)
\]
相关，分布相同时达到最优。这个结论依赖判别器足够强、训练达到平衡以及分布可比较。有限数据和不充分优化会破坏这些条件。

\subsection{逼真与覆盖为何会冲突}

GAN 不需要显式计算$p_g(x)$，因此可以把容量用于生成清晰样本。它的困难来自两个网络的交替博弈。判别器过强时，生成器获得的梯度可能饱和。不同噪声被映射到相似样本时，称为模式坍塌。生成结果可能很逼真，却遗漏罕见故障或某些工况。条件 GAN 可以把类别、工况或部分观测输入生成器和判别器，但条件一致性仍需单独验证。

GAN 适合采样速度要求高、且有足够样本支持目标分布的任务。它不提供像 VAE 那样直接的后验编码，也不直接给出可比较的显式似然。训练稳定性、模式覆盖和条件满足率应与样本质量一起报告。

\section{扩散模型的逐步去噪}

\subsection{从一次生成改为多步恢复}

扩散模型把困难的生成映射分成许多局部去噪任务。前向过程逐步破坏真实样本，反向网络学习在每个噪声尺度恢复结构。设$x_0\sim p_{\mathrm{data}}$，$t=1,\ldots,T$，并令$0<\beta_t<1$：
\begin{equation}
q(x_t\mid x_{t-1})=\mathcal N(\sqrt{1-\beta_t}x_{t-1},\beta_tI).
\end{equation}
记$\alpha_t=1-\beta_t$、$\bar\alpha_t=\prod_{s=1}^t\alpha_s$，可直接得到
\begin{equation}
x_t=\sqrt{\bar\alpha_t}x_0+\sqrt{1-\bar\alpha_t}\epsilon,
\qquad \epsilon\sim\mathcal N(0,I).
\label{eq:generative-direct-noising}
\end{equation}
这里$t$是人为设定的噪声步，不是设备运行时间。训练时随机选$t$和$\epsilon$，所以不必真的执行完整的前向链。

例如$x_0=2$、$\bar\alpha_t=0.64$且$\epsilon=1$时，$x_t=0.8\times2+0.6\times1=2.2$。网络看到$x_t$和$t$，不知道$x_0$。它要学习的是在给定噪声状态下最有用的恢复方向。

\subsection{噪声预测与反向采样}

噪声预测网络记为$\epsilon_{\theta}(x_t,t)$，常用训练目标为
\begin{equation}
\mathcal L_{\mathrm{simple}}=\E_{t,x_0,\epsilon}\left[\|\epsilon-\epsilon_{\theta}(x_t,t)\|_2^2\right].
\end{equation}
平方损失的最优函数是条件均值：
\[
\epsilon_{\theta}^*(x_t,t)=\E[\epsilon\mid x_t,t].
\]
由式\eqref{eq:generative-direct-noising}，可以得到干净样本估计
\[
\widehat x_0=\frac{x_t-\sqrt{1-\bar\alpha_t}\epsilon_{\theta}(x_t,t)}{\sqrt{\bar\alpha_t}}.
\]
当$\bar\alpha_t$很小时，噪声预测误差会被放大。这说明高噪声阶段的干净样本反推对噪声误差敏感，但不能据此单独断定完整采样的不稳定程度。

采样从$x_T\sim\mathcal N(0,I)$开始，再按$t=T,T-1,\ldots,1$反向计算。以去噪扩散概率模型（Denoising Diffusion Probabilistic Model，DDPM）的一种参数化为例，
\begin{align}
\mu_{\theta}(x_t,t)&=\frac1{\sqrt{\alpha_t}}
\left(x_t-\frac{\beta_t}{\sqrt{1-\bar\alpha_t}}\epsilon_{\theta}(x_t,t)\right),\\
x_{t-1}&=\mu_{\theta}(x_t,t)+\sqrt{\widetilde\beta_t}\,\zeta_t,
\end{align}
其中$\zeta_t\sim\mathcal N(0,I)$是反向步骤重新抽取的噪声。通常在$t>1$时加入该随机项，在最后一步令其为零，以避免向最终样本额外加噪。该式采用固定的后验方差$\widetilde\beta_t=\beta_t(1-\bar\alpha_{t-1})/(1-\bar\alpha_t)$；其他方差参数化会改变采样分布。一次计算$\widehat x_0$不等于执行完整反向采样。

完整的变分目标包含前向后验与反向转移之间的 KL 项。常用的噪声均方误差是经过参数化和权重处理后的简化目标，不能与完整变分目标逐项等同。噪声预测、$x_0$预测和速度预测改变梯度尺度与采样表现，但共享逐步去噪框架。

\subsection{条件控制与扩散模型的变体}

给定条件$c$时，网络写成$\epsilon_{\theta}(x_t,t,c)$。无分类器引导让同一网络同时学习有条件和无条件预测，并组合为
\begin{equation}
\widehat\epsilon=\epsilon_{\theta}(x_t,t,\varnothing)+w\left[\epsilon_{\theta}(x_t,t,c)-\epsilon_{\theta}(x_t,t,\varnothing)\right].
\end{equation}
$w$越大，条件方向通常越强，但多样性可能下降并产生伪影。边缘、深度、掩码或几何信息也可以作为条件输入。条件输入不自动保证物理可行性。

去噪扩散隐式模型（Denoising Diffusion Implicit Model，DDIM）使用不同的反向轨迹减少采样步数。潜空间扩散先把数据压缩到潜空间，再在潜空间进行扩散，降低高分辨率任务的计算量。这些变体改变采样速度或计算位置，没有消除条件一致性和压缩损失的验证要求。

\section{任务选择与验证边界}

VAE、GAN 和扩散模型解决的是不同的计算折中。需要连续潜变量、补全或稳定编码时，VAE 的后验结构较直接。需要快速采样且主要关注视觉或波形逼真度时，可以考虑 GAN，但必须检查模式坍塌。需要复杂条件控制、较好的覆盖性和高质量样本时，扩散模型通常更合适，但采样成本较高。若任务要求显式似然、可逆映射或异常分数，还需要单独评估其他密度建模方法，不能从上述模型的采样质量推断得到。

工业应用至少要检查四类性质：
\begin{enumerate}
  \item \textbf{真实性：}边缘分布、相关性、频谱或缺陷几何是否接近真实数据。
  \item \textbf{覆盖性：}不同设备、工况和故障类型是否都出现，而不是重复少数模板。
  \item \textbf{条件与物理一致性：}生成结果是否符合输入条件、守恒关系、边界条件和状态转移约束。
  \item \textbf{下游价值：}加入合成数据后，独立真实测试集上的召回、误报和校准是否改善。
\end{enumerate}

约束可以加入生成损失，例如
\begin{equation}
\mathcal L=\mathcal L_{\mathrm{gen}}+\lambda_{\mathrm{phys}}\E[\|r(x,c)\|_2^2],
\end{equation}
其中$r$是物理残差。惩罚过弱会产生不可行样本，过强则可能损害分布覆盖。应分别报告真实性、多样性和约束违反率。

生成数据还可能记住训练集中的罕见记录，或改变真实故障的先验比例。测试集必须保持真实、独立，并且不能参与生成模型训练或选择。评价时应与“只用真实数据”的基线比较，匹配训练预算，并在不同随机种子、设备和故障类型上重复。生成样本的数量增加、训练损失下降或单张图片更清晰，都不能单独证明下游收益。

\section*{辅助学习材料}
\begingroup
\emergencystretch=3em
\begin{itemize}
\item \textbf{What are Diffusion Models?}；作者/平台：Lilian Weng；英文技术教程。可对照前向加噪、反向去噪、条件引导和潜空间扩散。\href{https://lilianweng.github.io/posts/2021-07-11-diffusion-models/}{在线阅读}。
\item \textbf{From Autoencoder to Beta-VAE}；作者/平台：Lilian Weng；英文技术教程。可对照 VAE、ELBO 和重参数化。\href{https://lilianweng.github.io/posts/2018-08-12-vae/}{在线阅读}。
\item \textbf{第六节 2021--生成式对抗网络}；作者/平台：李宏毅；中文课程。可对照 GAN 的生成器、判别器和训练过程。\href{https://www.bilibili.com/video/BV1Wv411h7kN/?p=58}{观看分集}。
\end{itemize}
\endgroup

\paragraph{本章小结。}
生成任务首先要明确输出分布和条件。VAE 通过潜变量和 ELBO 获得可采样表示，GAN 通过判别器追求逼真但容易训练不稳定和模式坍塌，扩散模型通过多步去噪处理复杂分布并提供条件控制。模型选择不能只看样本外观。真实性、覆盖性、条件与物理一致性，以及独立真实数据上的下游结果共同决定生成模型是否适用。